\documentclass[11pt]{article}
\usepackage{amsmath,amssymb,amsthm}
\usepackage{booktabs}
\usepackage[margin=1in]{geometry}
\newtheorem{theorem}{Theorem}
\newtheorem{lemma}[theorem]{Lemma}
\title{Problem 786: Billiard}
\author{Project Euler Solutions}
\date{}
\begin{document}
\maketitle
\section{Problem Statement}
A billiard table is a quadrilateral with angles $120^\circ, 90^\circ, 60^\circ, 90^\circ$ at vertices $A, B, C, D$ respectively, with $AB = AD$. A ball departs from $A$, bounces elastically off edges (never at corners), and returns to $A$. Let $B(N)$ be the number of possible traces with at most $N$ bounces. Given $B(10)=6$, $B(100)=478$, $B(1000)=45790$. Find $B(10^9)$.
\section{Analysis}
\subsection*{Unfolding the Billiard Table}

The classical technique for analyzing billiard trajectories is **unfolding**: instead of reflecting the ball's path at each bounce, we reflect the table and follow the straight-line trajectory in the unfolded plane.

\subsection*{Tiling Property}

The quadrilateral with angles $120^\circ + 90^\circ + 60^\circ + 90^\circ = 360^\circ$ tiles the plane perfectly. Reflecting across each edge produces a periodic tiling. Therefore, the unfolded trajectory is a straight line from the origin to a copy of vertex $A$ in the tiling lattice.

\subsection*{Lattice Structure}

The tiling lattice has basis vectors deter
\section{Answer}

\[
\boxed{45594532839912702}
\]

\end{document}
